\chapter{Conclusion}
\label{chap:conclusion}
...
% \chapter{Перспективные направления и дальнейшее развитие работы}
% \label{chap:future_work}

% Результаты, полученные в ходе данного исследования, подтвердили высокую эффективность предложенных моделей (Phy-SRGAN и ST-PhyGAN) для решения задачи Super-Resolution динамики многофазных течений. Тем не менее, область применения глубокого обучения для ускорения гидродинамических расчетов активно развивается, и существует ряд перспективных направлений, которые могут лечь в основу дальнейшего развития данной работы. В этом разделе будут рассмотрены две гипотезы, основанные на самых передовых архитектурных и методологических подходах в современном машинном обучении.

% \section{Повышение точности за счет глобального контекста: архитектуры на основе Vision Transformers}
% \label{sec:future_vit}

% \textbf{Ограничение текущего подхода.} Модели, основанные на сверточных (CNN) и сверточно-рекуррентных (ConvLSTM) слоях, такие как наш ST-PhyGAN, превосходно улавливают локальные пространственные и временные зависимости. Их работа основана на принципе "локального рецептивного поля": каждый нейрон анализирует лишь небольшую окрестность. Однако гидродинамические процессы в пласте по своей природе являются глобальными: изменение давления в добывающей скважине в одной части месторождения мгновенно влияет на поле скоростей во всех остальных частях. Существующие сверточные архитектуры с трудом улавливают такие дальнодействующие физические корреляции.

% \textbf{Предлагаемая гипотеза: Spatio-Temporal Vision Transformer GAN (ST-ViT-GAN).} Мы предполагаем, что дальнейшее повышение точности и физического реализма предсказаний может быть достигнуто за счет интеграции \textbf{механизма внимания (self-attention)}, лежащего в основе архитектуры \textbf{Vision Transformer (ViT)}. Предлагается модифицировать архитектуру генератора ST-PhyGAN, заменив сверточные слои в его наиболее сжатой части ("бутылочном горлышке") на блок трансформера.

% Механизм внимания позволит модели для каждой точки пространства вычислять ее взаимосвязь со всеми остальными точками, эффективно создавая глобальное рецептивное поле. Это даст сети возможность напрямую моделировать дальнодействующие зависимости в поле давления, что, в свою очередь, должно привести к более точному предсказанию траекторий зарождения и роста "вязких пальцев", зависящих от общей картины течения.

% \textbf{Новизна.} Применение гибридных CNN-Transformer архитектур для пространственно-временного супер-разрешения нестабильных многофазных течений является абсолютно новым подходом, который потенциально может стать следующим шагом в развитии суррогатных моделей для задач гидродинамики.

% \section{Радикальное ускорение: дистилляция знаний для интерактивных расчетов}
% \label{sec:future_kd}

% \textbf{Ограничение текущего подхода.} Наша финальная модель ST-PhyGAN, несмотря на на порядки большее быстродействие по сравнению с классическим симулятором, все еще может быть достаточно "тяжелой" с точки зрения вычислительных ресурсов. Это может ограничивать ее применение в задачах, требующих практически мгновенного отклика, например, в интерактивных инженерных тренажерах или системах поддержки принятия решений в реальном времени.

% \textbf{Предлагаемая гипотеза: Дистилляция знаний (Knowledge Distillation) для создания сверхлегких моделей.} Мы предполагаем, что возможно создать сверхбыструю суррогатную модель практически без потери качества путем применения методологии дистилляции знаний. Процесс состоит из двух этапов:
% \begin{enumerate}
%     \item \textbf{Обучение модели-"учителя":} Сначала обучается наша наиболее сложная и точная модель, например, ST-PhyGAN. Эта модель выступает в роли "эксперта".
%     \item \textbf{Обучение модели-"студента":} Затем создается очень компактная и быстрая нейросеть (например, небольшой U-Net без рекуррентных и состязательных компонентов). Эта модель-"студент" обучается не на "сырых" данных из симулятора, а на задаче имитации выходных данных модели-"учителя".
% \end{enumerate}

% Таким образом, "учитель" передает "студенту" свои обобщенные "знания" о процессе. В результате "студент" учится аппроксимировать сложное поведение "учителя", но делает это на порядки быстрее за счет меньшего количества параметров.

% \textbf{Новизна.} Применение дистилляции знаний для создания "суррогатной модели от суррогатной модели" с целью достижения производительности, достаточной для интерактивных расчетов в задачах гидродинамического моделирования, является новой и крайне практико-ориентированной идеей. Успешная реализация этого подхода может сделать продвинутое моделирование доступным не только на высокопроизводительных кластерах, но и на персональных рабочих станциях инженеров.